\subsection{冻结截断、容量全息与统计审计的三寄存器分裂}\label{subsec:conclusion-freezing-capacity-samplepattern-audit-register-trisplitting}
沿用小节 \ref{subsec:conclusion-maxfiber-freezing-capacity-entropy-separation}、小节 \ref{subsec:conclusion-capacity-mellin-thermodynamic-wedderburn-holography}、小节 \ref{subsec:conclusion-capacity-majorization-infonce-complexity-separation} 与定理 \ref{thm:conclusion-sample-pattern-family-closes-hot-thermodynamics} 的记号。前述结果已经分别给出：亚极大预算下仅最大纤维受截断的容量平台、冻结相 escort 测度在全 Rényi 轴与全变差上的指数塌缩、连续容量曲线与 Wedderburn 块谱的完备对应、有限样本相等模式对对称接口的全恢复，以及完整 Bayes--InfoNCE 数值塔对有序纤维谱的最小连续层析。把这些接口并读之后，还可抽出一条更高层的闭合链：冻结前沿、容量全息、样本模式热相几何与统计审计塔并非同一信息的不同近似，而是三条彼此不可替代的 faithful 坐标轴。

\begin{theorem}[亚极大预算窗的单层截断律]\label{thm:conclusion-subcritical-budget-single-layer-truncation}
\leanverified{paper\_conclusion\_subcritical\_budget\_single\_layer\_truncation}
记
\[
M_m:=\max_{x\in X_m}d_m(x),
\qquad
K_m:=\#\{x\in X_m:\ d_m(x)=M_m\},
\]
\[
M_m^{(2)}:=\max\{d_m(x):\ d_m(x)<M_m\},
\qquad
\alpha_\ast:=\lim_{m\to\infty}\frac{1}{m}\log M_m,
\]
\[
g_\ast:=\lim_{m\to\infty}\frac{1}{m}\log K_m,
\qquad
\alpha_\ast^{(2)}:=\limsup_{m\to\infty}\frac{1}{m}\log M_m^{(2)}.
\]
设整数预算序列 \((B_m)_{m\ge 1}\) 满足
\[
\beta:=\lim_{m\to\infty}\frac{1}{m}\log 2^{B_m}
\]
存在，且
\[
\alpha_\ast^{(2)}<\beta<\alpha_\ast.
\]
则对充分大的 \(m\) 恰有
\[
C_m(B_m)=2^m-K_m\bigl(M_m-2^{B_m}\bigr),
\]
从而
\[
1-\mathrm{Succ}_m(B_m)
=
\frac{K_m(M_m-2^{B_m})}{2^m}.
\]
特别地，
\[
\lim_{m\to\infty}\frac{1}{m}\log\bigl(1-\mathrm{Succ}_m(B_m)\bigr)
=
\alpha_\ast+g_\ast-\log 2.
\]
因此，在该预算窗内，预算 \(2^{B_m}\) 只截断最大纤维层，而对一切严格次大纤维完全透明。
\end{theorem}
\begin{proof}
定理 \ref{thm:conclusion-subcritical-maxfiber-capacity-platform} 已给出
\[
C_m(B_m)=2^m-K_m\bigl(M_m-2^{B_m}\bigr)
\]
及其指数率结论。其证明同时表明：由
\[
\alpha_\ast^{(2)}<\beta<\alpha_\ast
\]
可得对充分大的 \(m\) 有
\[
M_m^{(2)}<2^{B_m}<M_m,
\]
故所有严格次大纤维都被预算完全覆盖，而仅最大纤维被截断。再把
\[
\mathrm{Succ}_m(B)=\frac{C_m(B)}{2^m}
\]
代入即可得到
\[
1-\mathrm{Succ}_m(B_m)
=
\frac{K_m(M_m-2^{B_m})}{2^m},
\]
从而得到所述单层截断律。证毕。
\end{proof}

\begin{theorem}[冻结相的全轴熵与全变差塌缩]\label{thm:conclusion-freezing-measure-collapse-all-axes}
\leanverified{paper\_conclusion\_freezing\_measure\_collapse\_all\_axes}
定义冻结阈值
\[
a_{\mathrm c}:=\frac{\log\varphi-g_\ast}{\alpha_\ast-\alpha_\ast^{(2)}},
\]
并在 \(a>a_{\mathrm c}\) 时定义 escort 分布
\[
\pi_{m,a}(x):=\frac{d_m(x)^a}{S_a(m)},
\qquad
S_a(m):=\sum_{x\in X_m}d_m(x)^a.
\]
再记
\[
A_m^{\max}:=\{x\in X_m:\ d_m(x)=M_m\},
\qquad
U_m^{\max}:=\mathrm{Unif}(A_m^{\max}).
\]
则对每个 \(a>a_{\mathrm c}\) 都有
\[
\pi_{m,a}(A_m^{\max})
\ge
1-\exp\!\bigl(-m\Delta(a)+o(m)\bigr),
\]
\[
\Delta(a):=a(\alpha_\ast-\alpha_\ast^{(2)})-(\log\varphi-g_\ast)>0,
\]
并且对每个 \(s\in[1,\infty]\),
\[
\lim_{m\to\infty}\frac{1}{m}H_s(\pi_{m,a})=g_\ast.
\]
此外还有精确恒等式
\[
\bigl|\pi_{m,a}-U_m^{\max}\bigr|_{\mathrm{TV}}
=
1-\pi_{m,a}(A_m^{\max})
\le
\exp\!\bigl(-m\Delta(a)+o(m)\bigr).
\]
因此，冻结相并不只是压力函数在高阶矩上的线性化，而是 escort 测度在质量、Shannon/R\'enyi 熵率与全变差几何上的同步塌缩。
\end{theorem}
\begin{proof}
由定理 \ref{thm:conclusion-frozen-escort-shannon-renyi-collapse}，当 \(a>a_{\mathrm c}\) 时有
\[
\pi_{m,a}(A_m^{\max})
\ge
1-\exp\!\bigl(-m\Delta(a)+o(m)\bigr),
\]
并且 Shannon 熵率与任意固定有限 Rényi 阶都收敛到 \(g_\ast\)。再由推论 \ref{cor:conclusion-frozen-escort-all-renyi-collapse}，该结论延伸到全部 \(s\in[1,\infty]\)。最后，定理 \ref{thm:conclusion-frozen-escort-tv-rigidity} 给出精确全变差恒等式
\[
\bigl|\pi_{m,a}-U_m^{\max}\bigr|_{\mathrm{TV}}
=
1-\pi_{m,a}(A_m^{\max}),
\]
以及相同指数率的上界。三者合并即得结论。证毕。
\end{proof}

\begin{theorem}[连续容量曲线的热力学--Wedderburn 完备性]\label{thm:conclusion-capacity-wedderburn-thermodynamic-completeness}
\leanverified{paper\_conclusion\_capacity\_wedderburn\_thermodynamic\_completeness}
固定分辨率 \(m\)。定义连续容量曲线
\[
\mathcal C_m^{\mathrm{cont}}(T):=\sum_{x\in X_m}\min(d_m(x),T)
\qquad (T\ge 0),
\]
尾计数函数
\[
N_m(T):=\#\{x\in X_m:\ d_m(x)\ge T\},
\]
直方图
\[
h_m(d):=\#\{x\in X_m:\ d_m(x)=d\},
\]
以及纤维群胚代数
\[
A_m:=\mathbf C[R_m^{\mathrm{bin}}].
\]
则下列四类对象彼此唯一决定：
\[
\mathcal C_m^{\mathrm{cont}}(\cdot)
\Longleftrightarrow
N_m(\cdot)
\Longleftrightarrow
h_m(\cdot)
\Longleftrightarrow
A_m\simeq \bigoplus_{d\ge 1}M_d(\mathbf C)^{\oplus h_m(d)}.
\]
并且对任意 \(q>0\)，有
\[
S_q(m):=\sum_{x\in X_m}d_m(x)^q
=
q\int_0^\infty T^{q-1}N_m(T)\,\dd T,
\]
从而全部正阶矩谱也由 \(\mathcal C_m^{\mathrm{cont}}\) 唯一恢复。

更强地，若
\[
f:\Omega\to X,
\qquad
g:X\to Y,
\qquad
h:=g\circ f,
\]
并记
\[
X_+:=\{x\in X:\ d_f(x)>0\},
\]
则
\[
\mathcal C_h^{\mathrm{cont}}(T)=\mathcal C_f^{\mathrm{cont}}(T)
\qquad (\forall T\ge 0)
\]
当且仅当 \(g|_{X_+}\) 单射。因而，任何真正合并正纤维的粗粒化，都不可能同时在容量曲线、全部正阶矩谱与 Wedderburn 块谱三侧隐身。
\end{theorem}
\begin{proof}
定理 \ref{thm:conclusion-wedderburn-threshold-holography} 已证明
\[
\mathcal C_m^{\mathrm{cont}}(\cdot)
\Longleftrightarrow
N_m(\cdot)
\Longleftrightarrow
h_m(\cdot)
\Longleftrightarrow
A_m\simeq \bigoplus_{d\ge 1}M_d(\mathbf C)^{\oplus h_m(d)}.
\]
另一方面，对任意 \(q>0\) 与任意正整数 \(d\)，都有
\[
d^q=q\int_0^d T^{q-1}\,\dd T.
\]
对 \(d=d_m(x)\) 求和并交换求和与积分，得到
\[
S_q(m)
=
q\int_0^\infty T^{q-1}\#\{x\in X_m:\ d_m(x)\ge T\}\,\dd T
=
q\int_0^\infty T^{q-1}N_m(T)\,\dd T.
\]
故全部正阶矩谱也由 \(N_m\)，从而由 \(\mathcal C_m^{\mathrm{cont}}\) 唯一决定。

最后，粗粒化判据正是定理 \ref{thm:conclusion-coarsening-global-rigidity-of-capacity-curve} 的结论：若 \(g|_{X_+}\) 非单射，则存在两个正纤维被真合并，导致某个阈值区间上
\[
\mathcal C_h^{\mathrm{cont}}(T)<\mathcal C_f^{\mathrm{cont}}(T);
\]
反之若 \(g|_{X_+}\) 单射，则正纤维多重集只被重标号，容量曲线保持不变。证毕。
\end{proof}

\begin{theorem}[样本相等模式与 Bayes 深尾律的双参数闭包]\label{thm:conclusion-sample-pattern-deeptail-top-geometry-completeness}
固定分辨率 \(m\)，记
\[
N:=|X_m|,
\qquad
w_i:=\frac{d_i}{2^m}
\qquad (1\le i\le N)
\]
为按非增序排列的纤维权重。令 \(\Pi_N\) 表示 \(N\) 个独立样本
\[
X_1,\dots,X_N\sim (w_1,\dots,w_N)
\]
的相等模式分割，并记 Pareto 上边界为 \(H_m^{\mathrm{orc}}:=H_m\)。则 \(\mathrm{Law}(\Pi_N)\) 唯一决定
\[
\{w_i\}_{i=1}^{N},
\qquad
\{S_q(m)\}_{q\in\RR},
\qquad
H_m^{\mathrm{orc}},
\qquad
\{\mathrm{Succ}^{\mathrm{orc}}_m(B)\}_{B\ge 0}.
\]

进一步，若
\[
w_1=\cdots=w_{r_\ast}=w_\ast>w_{r_\ast+1}\ge\cdots\ge w_N>0,
\]
并定义
\[
h_t(w):=[z^t]\prod_{i=1}^{N}(1-w_i z)^{-1},
\qquad
T(\varepsilon):=\min\{t\ge 0:\ h_t(w)\le \varepsilon\},
\]
则存在常数 \(C_w>0\) 使
\[
h_t(w)=C_w\,t^{r_\ast-1}w_\ast^t\bigl(1+O(t^{-1})\bigr),
\]
从而
\[
w_\ast=\lim_{t\to\infty}\frac{h_{t+1}(w)}{h_t(w)},
\qquad
r_\ast
=
1+\lim_{t\to\infty}
t\left(\frac{h_{t+1}(w)}{w_\ast h_t(w)}-1\right),
\]
并且当 \(\varepsilon\downarrow 0\) 时，
\[
T(\varepsilon)
=
\frac{\log(1/\varepsilon)}{\log(1/w_\ast)}
+
\frac{r_\ast-1}{\log(1/w_\ast)}\log\log(1/\varepsilon)
+
O(1).
\]
因此，有限样本相等模式恢复的是完整有限谱，而深尾 Bayes 查询恢复的是顶端权 \(w_\ast\) 与顶端重数 \(r_\ast\) 的双参数几何。
\end{theorem}
\begin{proof}
定理 \ref{thm:conclusion-sample-pattern-complete-symmetric-interface} 已经说明：\(\mathrm{Law}(\Pi_N)\) 唯一决定有序权重多重集、全部幂和谱、Pareto 上边界以及整条有限分辨率预言机成功曲线。另一方面，定理 \ref{thm:conclusion-bayes-full-recovery-top-pole} 给出
\[
h_t(w)=C_w\,t^{r_\ast-1}w_\ast^t\bigl(1+O(t^{-1})\bigr),
\]
并由相邻比值恢复 \(w_\ast\) 与 \(r_\ast\)。最后，推论 \ref{cor:conclusion-bayes-deep-tail-two-parameter-log-law} 正给出阈值
\[
T(\varepsilon)
\]
的双参数对数律。三者合并即得结论。证毕。
\end{proof}

\begin{theorem}[热相接触几何由样本模式族完全闭合]\label{thm:conclusion-sample-pattern-family-closes-hot-contact-geometry}
设热相区间为 \((a_0,a_1)\)，记预言机容量指数为 \(\Gamma_{\mathrm{orc}}(a)\)，共轭参数为 \(q_\ast(a)\in(0,1)\)，实参数压力为 \(\tau(q)\)，多重分形前沿为 \(f(a)\)。则在热相内有严格互反公式
\[
q_\ast(a)=1-\Gamma_{\mathrm{orc}}'(a),
\qquad
a=\tau'(q_\ast(a)),
\]
\[
\tau(q_\ast(a))
=
\Gamma_{\mathrm{orc}}(a)-a\,\Gamma_{\mathrm{orc}}'(a),
\qquad
\tau''(q_\ast(a))=-\frac{1}{\Gamma_{\mathrm{orc}}''(a)}.
\]
并且
\[
\Gamma_{\mathrm{orc}}(a)=f(a)+a
\]
在整个热相中没有幽灵扇区：每个 \(a\in(a_0,a_1)\) 都由唯一真实接触点 \(\alpha=a\) 暴露。更强地，跨尺度样本模式族
\[
\{\mathrm{Law}(\Pi_{|X_m|})\}_{m\ge 1}
\]
唯一决定热相中的全部对象
\[
f(a),
\qquad
\tau(q),
\qquad
a\longleftrightarrow q_\ast(a).
\]
\end{theorem}
\begin{proof}
定理 \ref{thm:conclusion-hot-phase-oracle-legendre-reciprocity} 已给出题述四个互反公式。定理 \ref{thm:conclusion-hot-phase-no-ghost-sector} 进一步说明
\[
\Gamma_{\mathrm{orc}}''(a)<0,
\qquad
a\longmapsto q_\ast(a)
\]
是从 \((a_0,a_1)\) 到 \((0,1)\) 的 \(C^1\)-微分同胚，故热相内不存在伪暴露扇区。最后，定理 \ref{thm:conclusion-sample-pattern-family-closes-hot-thermodynamics} 已证明跨尺度样本模式族唯一决定 \(\Gamma_{\mathrm{orc}}\)，从而经由上述互反公式唯一恢复 \(f\)、\(\tau\) 与 \(a\leftrightarrow q_\ast(a)\)。证毕。
\end{proof}

\begin{theorem}[主序--容量--完整 Bayes--InfoNCE 塔的严格反向同构]\label{thm:conclusion-majorization-capacity-infonce-strict-antiisomorphism}
\leanverified{paper\_conclusion\_majorization\_capacity\_infonce\_strict\_antiisomorphism}
固定总质量 \(S\) 与可见类型数 \(N\)，记
\[
\mathcal S_{S,N}
:=
\left\{
\mathbf d=(d_1,\dots,d_N)\in \ZZ_{\ge 1}^N:
d_1\ge\cdots\ge d_N,\ \sum_{i=1}^{N}d_i=S
\right\}.
\]
对每个 \(\mathbf d\in\mathcal S_{S,N}\)，定义连续容量曲线
\[
C_{\mathbf d}(u):=\sum_{i=1}^{N}\min(d_i,u)
\qquad (u\ge 0),
\]
并令
\[
\mathscr L_{\mathbf d}
:=
\bigl(L_2^\star(\mathbf d),L_3^\star(\mathbf d),\dots,L_N^\star(\mathbf d)\bigr)
\]
表示把归一化谱 \(\mathbf d/S\) 代入 fixed-resolution Bayes--InfoNCE 最优闭式后所得的完整数值塔。则
\[
\mathbf d\succcurlyeq_{\mathrm{HLP}}\mathbf e
\iff
C_{\mathbf d}(u)\le C_{\mathbf e}(u)\ \ (\forall u\ge 0)
\iff
\mathscr L_{\mathbf d}\preccurlyeq_{\mathrm{InfoNCE}}\mathscr L_{\mathbf e},
\]
其中 \(\succcurlyeq_{\mathrm{HLP}}\) 为 Hardy--Littlewood--P\'olya 主序，而右端偏序由 \(\mathbf d\mapsto \mathscr L_{\mathbf d}\) 唯一搬运得到。

特别地，在 window-\(6\) 上同时成立
\[
B_{\mathrm{full}}^{\mathrm{bin}}(6)=2,
\]
\[
\text{不存在任何有限阿贝尔群上的线性陪集划分，其纤维划分等于 }\Fold_6^{\mathrm{bin}}\text{ 的纤维划分},
\]
\[
\text{在 }U_{21}\subset\RR^{21}\text{ 上，少于 }20=|X_6|-1\text{ 个 Bayes 最优 InfoNCE 标量不可能连续单射恢复有序纤维谱}.
\]
因此，window-\(6\) 上至少并存三条互不塌缩的复杂度轴：实例反演复杂度、线性表示复杂度与统计审计复杂度。
\end{theorem}
\begin{proof}
前一组等价正是推论 \ref{cor:conclusion-capacity-majorization-infonce-antiisomorphism} 的内容：在固定 \((S,N)\)-切片上，主序、容量曲线逐点序与完整 Bayes--InfoNCE 塔上的搬运偏序严格一致。window-\(6\) 的三条特化则由推论 \ref{cor:conclusion-window6-inversion-linear-statistical-triple-separation} 直接给出。证毕。
\end{proof}

\begin{conjecture}[freeze--capacity--audit 支链上的 faithful \texorpdfstring{$\pi$}{pi} 公式三寄存器下界]\label{conj:conclusion-piformula-freezing-capacity-audit-three-register-lower-bound}
在猜想 \ref{conj:conclusion-foldgauge-pi-three-register-lower-bound} 的语义下，若进一步要求 projective \(\pi\)-通道 faithful 地承载本支链中的冻结前沿、容量全息与统计审计结构，则其外置寄存器至少应裂为如下三类彼此正交的层：
\[
\mathfrak R_{\max},
\qquad
\mathfrak R_{\mathrm{cap}},
\qquad
\mathfrak R_{\mathrm{aud}}.
\]
其中
\begin{enumerate}
  \item \(\mathfrak R_{\max}\) 负责记录
  \[
  (M_m,K_m,M_m^{(2)}),
  \qquad
  (\alpha_\ast,g_\ast,\alpha_\ast^{(2)}),
  \]
  并据此支配定理 \ref{thm:conclusion-subcritical-budget-single-layer-truncation} 与定理 \ref{thm:conclusion-freezing-measure-collapse-all-axes} 所刻画的单层截断律与冻结相测度塌缩；
  \item \(\mathfrak R_{\mathrm{cap}}\) 负责记录
  \[
  \mathcal C_m^{\mathrm{cont}}(\cdot)
  \Longleftrightarrow
  N_m(\cdot)
  \Longleftrightarrow
  h_m(\cdot)
  \Longleftrightarrow
  A_m,
  \]
  并据此支配定理 \ref{thm:conclusion-capacity-wedderburn-thermodynamic-completeness} 中的热力学量、粗粒化可见性与 Wedderburn 块谱；
  \item \(\mathfrak R_{\mathrm{aud}}\) 负责记录
  \[
  \mathrm{Law}(\Pi_{|X_m|}),
  \qquad
  h_t(w),
  \qquad
  \{L_{K,m}^\star\}_{K=2}^{|X_m|},
  \]
  并据此支配定理 \ref{thm:conclusion-sample-pattern-deeptail-top-geometry-completeness}、定理 \ref{thm:conclusion-sample-pattern-family-closes-hot-contact-geometry} 与定理 \ref{thm:conclusion-majorization-capacity-infonce-strict-antiisomorphism} 中的热相接触几何、深尾顶端双参数与统计审计门槛。
\end{enumerate}
该猜想断言：\(\mathfrak R_{\max}\) 只见冻结前沿而不见完整 Wedderburn 块谱；\(\mathfrak R_{\mathrm{cap}}\) 给出容量与半单全息，却不自动替代深尾顶端双参数与完整审计塔；\(\mathfrak R_{\mathrm{aud}}\) 虽闭合热相与审计几何，却不压缩为单独的最大纤维寄存器。因而，任何试图把上述三层重新压成单一值层标量接口的 \(\pi\)-公式，都至多在数值层给出局部近似，而不能在本支链上成为 faithful 的 completion 证书。
\end{conjecture}

\paragraph{统一读法}
这条链把冻结、容量与审计的接口压成了三种不可互换的完备性。亚极大预算窗的容量缺口只看见最大纤维层；冻结相 escort 测度在全部 Rényi 轴与全变差上一起塌缩到最大层均匀分布；连续容量曲线则把正阶矩谱、尾计数、直方图与 Wedderburn 半单块账本严格锁成同一对象；样本相等模式既恢复有限分辨率上的完整对称接口，又经深尾极点与热相接触公式闭合出顶端几何和热相多重分形几何；而固定总质量切片上的主序、容量与完整 Bayes--InfoNCE 数值塔又构成严格反向同构。于是，在本支链上 faithful 的 \(\pi\)-宿主不能退化为单一标量模型，而必须同时保留冻结寄存器、容量寄存器与审计寄存器三条彼此正交的结构轴。

\endinput
